\begin{longtable}{>{\raggedright\arraybackslash}p{0.19\textwidth}>{\raggedright\arraybackslash}p{0.43\textwidth}>{\raggedright\arraybackslash}p{0.3\textwidth}}
\toprule
Result & Declarations (module) & On paper only\\
\midrule
\endhead
Lemma~\ref{lem:micro} & \texttt{tracking\_gain\_explicit\_root}, \texttt{tracking\_prior\_riccati\_equation} (\nolinkurl{Microfoundations.lean}); the third example's recursion, \texttt{p1\_proportional\_trade\_recursion} (\nolinkurl{Paper1StressIdentification.lean}) & Convergence of sample covariances. The third example's first-order condition is assumed in the checked statement, not derived\\
Theorem~\ref{thm:eqinv} (sufficiency) & \texttt{p1\_eqAggFlow\_kernel}, \texttt{p1\_eqAggFlow\_of\_kernel}, \texttt{p1\_kernel\_of\_eqAggFlow}, \texttt{p1\_eq\_market\_invariance}, \texttt{p1\_eq\_insider\_profit\_invariance}, \texttt{p1\_eq\_group\_level\_differs} (\nolinkurl{Paper1EquilibriumInvariance.lean}) & The unbounded-horizon moment statement; the passage from pathwise equality to equal equilibria\\
Proposition~\ref{prop:anylinear} & \texttt{p1\_anylinear\_flow\_kernel}, \texttt{p1\_anylinear\_iff}, \texttt{p1\_anylinear\_invariance} (\nolinkurl{Paper1CoreExtensions.lean}) & The passage from pathwise equality to equal equilibria\\
Theorem~\ref{thm:inventory} & \texttt{p1\_inventory\_foc}, \texttt{p1\_inventory\_innovation}, \texttt{p1\_inventory\_benchmark} (\nolinkurl{Paper1InventoryEquilibrium.lean}); \texttt{p1\_inventory\_small\_feedback} (\nolinkurl{Paper1InventoryIFT.lean}) & Part (iii); the version of part (iv) uniform over small group measures\\
Theorem~\ref{thm:eqreduced} & \texttt{p1\_eq\_realize\_D}, \texttt{p1\_eq\_realize\_dealer}, \texttt{p1\_eq\_scaled\_cohort} (\nolinkurl{Paper1EquilibriumReducedForm.lean}) & Parts (ii)--(iii), and part (iv) beyond the realization\\
Theorem~\ref{thm:kcohort} & \texttt{p1\_thm\_kcohort\_lower\_bound\_uniform} (\nolinkurl{Paper1KCohortUniform.lean}); \texttt{p1\_thm\_kcohort\_lower\_bound\_weights}, \texttt{p1\_thm\_kcohort\_lower\_bound\_W1} (\nolinkurl{Paper1KCohortEstimators.lean}) & The matching lower order of the response gap; the checked pair is a binomial moment-matched pair of the same order\\
Theorem~\ref{prop:esprit} & \texttt{p1\_thm\_prony\_upper\_bound} (\nolinkurl{Paper1PronyUpper.lean}) & ---\\
Theorem~\ref{thm:minphase} & \texttt{p1\_thm\_minphase\_zeros} (\nolinkurl{Paper1MinPhase.lean}); \texttt{p1\_thm\_minphase\_spectrum\_H}, \texttt{p1\_thm\_minphase\_spectrum\_merged} (\nolinkurl{Paper1SpectralH.lean}); model link \texttt{p1\_eigen\_iff\_bottom\_root}, \texttt{p1\_schur\_stable\_iff\_poles\_in\_disk}, \texttt{p1\_impulse\_generating} (\nolinkurl{Paper1ModelLink.lean}) & ---\\
Theorem~\ref{prop:passive} & \texttt{p1\_thm\_passive\_cost\_matched}, \texttt{p1\_thm\_passive\_cost\_gauss\_radau\_two} (\nolinkurl{Paper1PassiveMatched.lean}); \nolinkurl{p1_thm_passive_cost_gauss_radau_three} (\nolinkurl{Paper1GaussRadau3.lean}); \texttt{p1\_thm\_passive\_cost\_every\_K} (\nolinkurl{Paper1PassiveKCohort.lean}) & The divergence-rate statements; the stationary limit; Gauss--Radau exactness for general $K$ enters as a hypothesis\\
Proposition~\ref{prop:target} & \texttt{p1\_prop\_target\_two\_cohorts}, \texttt{p1\_prop\_target\_every\_K} (\nolinkurl{Paper1SupervisorySeparation.lean}) & The checked two-group pair is a different moment-matched pair from the Gauss and Gauss--Radau pair, with the same orders; the every-$K$ version assumes a dominant-group margin; the largest-weight comparison holds for every $K\ge 2$ by Proposition~\ref{prop:maxweight}\\
Proposition~\ref{prop:targetK} and Lemma~\ref{lem:everyK} & The counting step of Lemma~\ref{lem:everyK}(iii) only: \texttt{p1\_largest\_value\_not\_constant}, \texttt{p1\_largest\_weight\_not\_constant} (\nolinkurl{Paper1EveryKHalt.lean}) & Lemma~\ref{lem:everyK}(i)--(ii), the finiteness of the level sets, and all of Proposition~\ref{prop:targetK}\\
Proposition~\ref{prop:nearcrit} & Steps 2 and 3: the lower bound $|D_A|\ge c_3(\epsilon+|\omega|)$ (\texttt{p1\_nearcrit\_lower}), the uniform ratio bound (\texttt{p1\_ratio\_bound}), $|u-1|\le4\rho$ and $u-1-\log u\le2(u-1)^2$, and the rate bound $16\rho^2$ from Whittle's formula (\texttt{p1\_nearcrit\_kl\_bound}) (\nolinkurl{Paper1NearCritical.lean}) & Whittle's formula; Step 1 (the kernel-gap bound); the local properties of $D^*$ from primitivity and their uniformity in $\Delta\eta$\\
Theorem~\ref{thm:partial} & Part (i) in full on joint price and tagged-flow records: \texttt{p1\_partial\_lower\_bound\_joint}, \texttt{p1\_partial\_lower\_bound\_matched} (\nolinkurl{Paper1PartialDisclosureLower.lean}); algebraic steps of (ii), (iii): \nolinkurl{Paper1PartialDisclosure.lean}; the deterministic core of (ii), subtract-then-Prony with existence, uniqueness and the weight bound $L\delta$, $\delta=2^{2K'}\eta(\Delta\eta)^{-(2K'-1)}$: \texttt{p1\_partial\_upper\_deterministic} (\nolinkurl{Paper1PartialDisclosureUpper.lean}) & The $O_P(N^{-1/2})$ rates of the tagged estimates that feed $\eta$ (Monte Carlo checks in the text)\\
Proposition~\ref{prop:hiding} & (i) from sufficiency (\texttt{hiding\_sufficiency}); (ii) the cap (\texttt{hiding\_cap}); (iii) adoption and its eventual form (\texttt{hiding\_adopt}, \texttt{hiding\_eventually}) (\nolinkurl{PublicationRevision8.lean}) & That targeting depends on the record only through its law, which is Theorem~\ref{thm:eqinv}\\
Proposition~\ref{prop:tagged} & \texttt{p1\_tagged\_iv\_determinant\_bound} (\nolinkurl{Paper1TaggedIV.lean}); \texttt{p1\_iv\_boundary\_decomposition} (\nolinkurl{Paper1IVBoundary.lean}) & The stationary IV central limit theorem\\
Proposition~\ref{prop:hidegame} & Parts (ii) and (iii) of the game (\texttt{p1\_hiding\_detection\_bound}, \texttt{p1\_hiding\_tag\_no\_crowding}, \nolinkurl{Paper1TrackingGame.lean}); part (iv), group by group (\texttt{p1\_hiding\_nash\_bound}, \nolinkurl{Paper1NewsLawHiding.lean}) & Part (i) of the game\\

Proposition~\ref{prop:nearcritn} & The eigenvalue step: $n$ eigenvalues within $4\rho$ of one give divergence at most $16n\rho^2$ (\texttt{p1\_nearcrit\_finite\_sample}, \nolinkurl{Paper1OpenResolved.lean}) & The Toeplitz sandwich\\
Companion Theorem~\ref{thm:cascade} & Rounds below every fixed point and convergence to the least one, the contraction bound, the jump (\texttt{p1\_cascade\_iterate\_le\_fixed}, \texttt{p1\_cascade\_rounds\_converge}, \texttt{p1\_cascade\_contraction}, \texttt{p1\_cascade\_jump}, \nolinkurl{Paper1OpenResolved.lean}) & The derivative step of (iii)\\
Companion Proposition~\ref{prop:insider} & $\mathrm{Re}\,(1-w)^{-1}\ge(1+s)^{-1}$, the symbol bound, netting of public flow (\texttt{p1\_re\_inv\_one\_sub\_ge}, \texttt{p1\_no\_manipulation\_symbol}, \texttt{p1\_netting\_prefix}, \nolinkurl{Paper1OpenResolved.lean}) & The Toeplitz representation, (i), (iii), and the step to conditional expectations in (iv)\\
Companion Proposition~\ref{prop:deadband} & The bias of a bounded deviation through a stable closed loop (\texttt{p1\_deadband\_bias}, \nolinkurl{Paper1OpenResolved.lean}); the Prony step of Theorem~\ref{prop:esprit} & The reduction of group flows to the recursion\\
Proposition~\ref{prop:gap} & none & All of it (the Gauss--Radau error formula is classical); exact check in 60-digit arithmetic for $K\le30$\\
Companion Proposition~\ref{prop:reecascade} & Three equilibria from the intermediate value theorem (\texttt{p1\_cascade\_three\_equilibria}, \nolinkurl{Paper1OpenResolved.lean}) & The equilibrium condition, (ii), (iv)\\
Companion Proposition~\ref{prop:smoothn} & The sup-norm contraction step, through \texttt{p1\_deadband\_bias} & The expansion, the $\ell^2$ bound and the test bound\\
Companion Proposition~\ref{prop:stressrate} & none & All of it\\
Companion Proposition~\ref{prop:learners} & The deterministic pathwise step (\texttt{p1\_flow\_market\_invariance}) & The randomized induction, (iii)\\
Companion Proposition~\ref{prop:insidergame} & The symbol bound behind (i), through \texttt{p1\_no\_manipulation\_symbol} & Existence, (ii)--(iv)\\
Companion Proposition~\ref{prop:cascadedyn} & The geometric bound of (i) and the price path of (iii) (\texttt{p1\_cascade\_rounds\_geometric}, \texttt{p1\_slow\_capital\_fall}, \nolinkurl{Paper1OpenResolved2.lean}) & The limits in (i) and (ii), the equilibrium condition in (iii), (iv)\\
Companion Proposition~\ref{prop:hardrate} & none & All of it\\
Companion Proposition~\ref{prop:nlcost} & The closed-loop bias bound in (i), through \texttt{p1\_deadband\_bias} & The Taylor bound, (ii), (iii)\\
Proposition~\ref{prop:maxweight} & The two polynomial inequalities behind the central binomial bounds (\texttt{p1\_centralbinom\_q\_step}, \texttt{p1\_centralbinom\_r\_step}, \nolinkurl{Paper1OpenResolved2.lean}) & The Sonin-function bounds and the asymptotics; exact rational check for $K\ge62$, 30-digit check for $K\le61$\\
Theorem~\ref{thm:questions}, Proposition~\ref{prop:parametric} & For finite populations: the tap identity and weak duality, the upper half of (iv) (\texttt{p1\_question\_weak\_duality}); the affine estimate of (i) (\texttt{p1\_question\_truncation\_error}); part (iii) (\texttt{p1\_long\_memory\_impossible}); the parametric direction of the proposition, modulus at most $\delta A+2\bar WA\bar r^L/\sqrt{1-\bar r^2}$ (\texttt{p1\_parametric\_modulus}) (\nolinkurl{Paper1QuestionsCost.lean}) & The lower half of (iv) (Sion's minimax theorem and the total-variation lower bound); (ii) (duality for best uniform approximation); the converse of the proposition (weak compactness in $\ell^2$); measures that are not finitely supported\\
Proposition~\ref{prop:endogenous_kappa} & The displayed bound of part (ii), as an arithmetic step: \nolinkurl{p1_endogenous_contraction_sensitivity} (\nolinkurl{Paper1ClaimRepairs.lean}) & Part (i), and the fixed-point argument of part (ii)\\
Proposition~\ref{prop:strategic} & Two arithmetic steps: \nolinkurl{p1_camouflage_unprofitable}, \nolinkurl{p1_sybil_share_count_bound} (\nolinkurl{Paper1ClaimRepairs.lean}) & The determinant factor of part (i) and the economic arguments\\
Proposition~\ref{prop:choice} & none & All of it\\
Theorem~\ref{thm:merged} and Theorem~\ref{M-n:ident} & (i) for the model: the price response obeys the closed loop with the cohort kernel, deconvolution, identification of the merged measure, its converse and identification from finitely many noise-free observations (\nolinkurl{p1_impulse_closed_loop_recursion}, \nolinkurl{p1_thm_merged_model_identification}, \nolinkurl{p1_thm_merged_model_converse}, \nolinkurl{Paper1MergedModel.lean}; \nolinkurl{p1_thm_merged_deconvolution}, \nolinkurl{p1_thm_merged_measure_identified}, \nolinkurl{Paper1Identification.lean}; \nolinkurl{p1_thm_merged_model_finite_identification}, \nolinkurl{Paper1FiniteIdentification.lean}); the pole count $1+K$ (\nolinkurl{Paper1PoleCount.lean}); (ii) the lower bound $\sigma^2\le2K_*^2P\delta^4T$ for tests with both errors at most $1/4$ (\nolinkurl{p1_thm_merged_separation_power}, \nolinkurl{Paper1MergedSeparation.lean}) & Minimality of the realization from lowest terms (a standard result of linear systems)\\
Theorem~\ref{thm:minphase}, part (ii) for every $K$ & The spectrum determines the memory rates and the merged measure, with possibly different innovation variances (\nolinkurl{Paper1SpectralID.lean}) & ---\\
Theorem~\ref{thm:kcohort}, further steps & The kernel moment gap $|k_m-k'_m|\le2W\binom mn d^n(c+d)^{m-n}$ (\nolinkurl{p1_thm_kcohort_kernel_moment_gap}, \nolinkurl{Paper1MomentGap.lean}); the lower bound from a kernel gap and its $K$-group form (\nolinkurl{p1_thm_lower_bound_from_kernel_gap}, \nolinkurl{p1_thm_kcohort_lower_bound_formal}, \nolinkurl{Paper1LowerBound.lean}); every error level (\nolinkurl{Paper1GeneralLevel.lean}) & ---\\
Theorem~\ref{thm:discvalue} (upper bounds) & In reset experiments, all $K$ merged weights from tagged first-period flows with error probability $Kv/(N\kappa^2\varepsilon^2)$, free of the spacing, and each memory rate with error $Le$ (\nolinkurl{p1_thm_tagged_recovery_all_classes}, \nolinkurl{Paper1Tagged.lean}; \nolinkurl{p1_thm_tagged_rate_recovery}, \nolinkurl{Paper1TaggedRates.lean}) & The stationary-record version (Proposition~\ref{prop:tagged})\\
Corollary~\ref{cor:unknown_primitives} & (i) persistence acts as a memory-rate-zero group; (ii) identification of $(\kappa,a,M)$ (\nolinkurl{p1_cor_unknown_primitives_recursion}, \nolinkurl{p1_cor_unknown_primitives_identification}, \nolinkurl{Paper1UnknownPrimitives.lean}) & Part (iii), the estimator's rate\\
\bottomrule
\end{longtable}
